% ============================================================
% Capitolo 11 — Vincoli SEC e PERF
% ============================================================

\chapter{Vincoli SEC e PERF}
\label{ch:vincoli-sec-perf}

Ogni progetto software ha requisiti di sicurezza e performance. Il problema è che questi requisiti vengono discussi una volta --- e poi evaporano dalla pratica quotidiana. I vincoli SEC e PERF di AI-DLC risolvono questo: sono una libreria di requisiti riusabili, definiti una volta nel framework, attivati in \texttt{\_CONTEXT.md}, riletti dall'agente prima di ogni operazione critica.

\section{La libreria SEC (sicurezza)}
\label{sec:11-libreria-sec}

Nove vincoli predefiniti in \texttt{SEC\_CONSTRAINTS.md}:

\begin{table}[htbp]
  \centering
  \begin{tabular}{lll}
    \toprule
    \textbf{ID} & \textbf{Nome} & \textbf{Principio} \\
    \midrule
    SEC-01 & Input Validation & Allow-list, validazione lato server obbligatoria \\
    SEC-02 & Authentication & Provider collaudato, token short-lived, no log secrets \\
    SEC-03 & Authorization & Least privilege, validazione ownership \\
    SEC-04 & Secrets \& Config & No hardcoded secrets, env vars, rotazione chiavi \\
    SEC-05 & Data Protection & PII minimizzata, redaction log, TLS \\
    SEC-06 & Dependency Hygiene & Pin versioni, SBOM, vulnerability scan \\
    SEC-07 & Threat Modeling & Asset, attori, STRIDE, rischio residuo (Fase 0--2) \\
    SEC-08 & Logging \& Monitoring & Log eventi sicurezza, strutturati JSON, no PII \\
    SEC-09 & SSRF & Allow-list URL in uscita, blocco IP privati \\
    \bottomrule
  \end{tabular}
  \caption{Libreria vincoli SEC}
  \label{tab:11-sec}
\end{table}

\section{La libreria PERF (performance)}
\label{sec:11-libreria-perf}

Sette vincoli predefiniti in \texttt{PERF\_CONSTRAINTS.md}:

\begin{table}[htbp]
  \centering
  \begin{tabular}{lll}
    \toprule
    \textbf{ID} & \textbf{Nome} & \textbf{Principio} \\
    \midrule
    PERF-01 & Latency Targets & SLO per endpoint (P50/P95/P99) \\
    PERF-02 & Throughput \& Concurrency & RPS target, async I/O, no blocking \\
    PERF-03 & Database Efficiency & Indici, no N+1, bound result sizes \\
    PERF-04 & Resource Utilization & Limiti CPU/memoria, profiling \\
    PERF-05 & Caching \& CDN & Cache layer, TTL, invalidazione \\
    PERF-06 & Startup \& Warm-up & Lazy loading, health check readiness \\
    PERF-07 & Build \& Bundle Size & Code splitting, tree shaking (frontend) \\
    \bottomrule
  \end{tabular}
  \caption{Libreria vincoli PERF}
  \label{tab:11-perf}
\end{table}

\section{Come attivare i vincoli in \texttt{\_CONTEXT.md}}
\label{sec:11-attivazione}

Non tutti i vincoli si applicano a ogni progetto. Attivi solo quelli rilevanti specificando la configurazione specifica del tuo contesto:

\begin{lstlisting}[language=,caption={Vincoli attivati per TaskFlow API}]
## Security Constraints
| SEC-01 | Input Validation  | Zod schema obbligatorio                    |
| SEC-02 | Authentication    | OIDC + JWT short-lived, refresh rotation   |
| SEC-03 | Authorization     | task.userId === auth.userId                |
| SEC-05 | Data Protection   | Pino redaction: password, token, email     |

## Performance Constraints
| PERF-01 | Latency     | P95 < 200ms — monitored via Grafana        |
| PERF-03 | DB Efficiency | Max 3 query per request, take: 50         |
\end{lstlisting}

La colonna \texttt{Spec} è la parte più importante: non ``usa autenticazione'' --- scrivi la specifica concreta che l'agente deve rispettare.

\section{Come l'agente usa i vincoli}
\label{sec:11-uso}

L'agente rilegge i vincoli attivi in due momenti:
\begin{itemize}
  \item Prima di generare codice in zone SEC-sensibili (autenticazione, input utente, logging, chiamate esterne)
  \item Prima di piani HIGH o superiori (il piano deve citare esplicitamente i vincoli rilevanti)
\end{itemize}

Il comando \texttt{@show-constraints} mostra in qualsiasi momento i vincoli attivi nella sessione corrente.

\section*{\LabelSummary}

\begin{itemize}
  \item SEC-01..09 coprono: input, autenticazione, autorizzazione, secrets, dati, dipendenze, threat modeling, logging, SSRF.
  \item PERF-01..07 coprono: latency, throughput, database, risorse, cache, startup, bundle.
  \item Si attivano in \texttt{\_CONTEXT.md} con la specifica concreta --- non il principio generico.
  \item L'agente li rilegge prima di codice SEC-sensibile e li cita nei piani HIGH+.
  \item Vincoli personalizzati si aggiungono con ID arbitrari (es. \texttt{SEC-10}, \texttt{PERF-08}).
\end{itemize}
